Fix \(P\in\mathcal P_{d,\epsilon,M,\sigma}\) and write
\[
 \vartheta:=\tau(P)/M,\qquad
 \vartheta_k:=\tau_k/M,\qquad
 s_{ak}:=P(X=k,A=a),\qquad
 z_{ak}:=E[(Y/M)\mathbf 1\{X=k,A=a\}].
\tag{1}
\]
Lemma~\(\mathrm{lem:scale\text{-}sanity}\) gives \(\vartheta\in[-1,1]\), and Assumption~\(\mathrm{ass:approximate\text{-}homogeneity}\) gives \(|\vartheta_k-\vartheta|\le\sigma\) on every positive-mass cell.

The three blocks in Definition~\(\mathrm{def:centered\text{-}product\text{-}estimator}\) are independent and each has size between \(n/4\) and \(n/2\) for \(n\ge4\).  Apply Lemma~\(\mathrm{lem:continuous\text{-}occupancy\text{-}collision\text{-}upper}\) to the pilot block.  After changing constants because its size is a fixed fraction of \(n\), the normalized pilot \(t\in[-1,1]\) satisfies
\[
 E[(t-\vartheta)^2]
 \le C_{\epsilon}\left\{\frac1n+\frac d{n^2}+\sigma^2\right\}.
\tag{2}
\]
The dimension condition for that lemma follows by taking the theorem's \(c_{\epsilon}\) small enough; all block sizes are deterministic.

Let \(\mathcal A\) be the classifier event on which every empirically light cell has \(p_k\le B/4\), while every empirically heavy cell has \(p_k\ge c_{0,\epsilon}L/n\).  We verify that this event is overwhelmingly likely.  If \(D\sim\operatorname{Binomial}(r,p)\), exponential Markov's inequality gives, after optimizing its positive or negative exponential parameter,
\[
 P(D\ge2rp)\le(e/4)^{rp},
 \qquad
 P(D\le rp/2)\le e^{-rp/8}.
\tag{3}
\]
For a cell with \(p_k>B/4\), its expected classifier occupancy is at least a fixed multiple of \(\beta_{\epsilon}L\), whereas the heavy threshold is \(\lambda_{\epsilon}L\).  For a cell with \(p_k<c_{0,\epsilon}L/n\), choose \(c_{0,\epsilon}\) so that its expected classifier occupancy is at most \(\lambda_{\epsilon}L/2\).  Applying \((3)\), using \(\beta_{\epsilon}=2^{12}\lambda_{\epsilon}\), and summing over at most
\[
 d\le \sqrt{\kappa_{\epsilon}}\,nL
\]
cells gives, after decreasing \(\kappa_{\epsilon}\) and increasing \(N_{\epsilon}\),
\[
 Q_P^{(n)}(\mathcal A^c)\le C_{\epsilon}n^{-8}.
\tag{4}
\]
The displayed numerical choice of \(\lambda_{\epsilon}\) also makes \(m\epsilon p_k\ge12L\) on every selected heavy cell.  Conditional on the classifier, independence of the estimation block therefore gives
\[
 P(N_{0k}N_{1k}=0)\le2e^{-m\epsilon p_k}\le2n^{-12}
\tag{5}
\]
for every heavy cell.

Condition on the pilot and classifier and suppose \(\mathcal A\) occurs.  For
\[
 r_{K,B}(s):=B^{-1}G_K(s/B),
 \qquad
 c_k:=s_{0k}s_{1k}\{r_{K,B}(s_{0k})+r_{K,B}(s_{1k})\},
\tag{6}
\]
the shifted-Chebyshev identity is
\[
 0\le x\{1-xG_K(x)\}\le K^{-2},
 \qquad 0\le x\le1.
\tag{7}
\]
Every light cell has \(p_k\le B/4\), and overlap makes \(s_{0k}\) and \(s_{1k}\) comparable.  Expanding
\[
 c_k-p_k
 =s_{1k}\{s_{0k}r_{K,B}(s_{0k})-1\}
 +s_{0k}\{s_{1k}r_{K,B}(s_{1k})-1\}
\]
and applying \((7)\) to the two terms gives
\[
 |c_k-p_k|\le C_{\epsilon}B/K^2.
\tag{8}
\]

For every partition \(\mathcal H\cup\mathcal L=\{1,\ldots,d\}\), the population centering identity is exact:
\[
 t+\sum_{k\in\mathcal H}p_k(\vartheta_k-t)
 +\sum_{k\in\mathcal L}c_k(\vartheta_k-t)-\vartheta
 =\sum_{k\in\mathcal L}(c_k-p_k)(\vartheta_k-t).
\tag{9}
\]
Since \(B\asymp_{\epsilon}L/n\), \(K\asymp L\), and
\[
 \max_{k:p_k>0}|\vartheta_k-t|^2
 \le2\{\sigma^2+(t-\vartheta)^2\},
\tag{10}
\]
equation \((8)\) and \(dB/K^2\le C_{\epsilon}d/(nL)=C_{\epsilon}\sqrt u\) show that the squared deterministic error in \((9)\) is at most
\[
 C_{\epsilon}u\{\sigma^2+(t-\vartheta)^2\}.
\tag{11}
\]

Lemma~\(\mathrm{lem:centered\text{-}mixed\text{-}arm\text{-}factorial\text{-}covariance}\), applied conditionally to the now deterministic light set, gives
\[
 \operatorname{Var}\!\left(
 \sum_{k\in\widehat{\mathcal L}}\widehat P_k(t)
 \mathrel{\Big|}t,I_{\mathrm{cls}},\mathcal A\right)
 \le C_{\epsilon}\left[
 \frac1n+\frac d{n^2}
 +u\{\sigma^2+(t-\vartheta)^2\}
 \right].
\tag{12}
\]
Its range condition holds because \(u\le\kappa_{\epsilon}\) implies \(d\le\sqrt{\kappa_{\epsilon}}\,nL\); choose \(\kappa_{\epsilon}\) so that this is at most \(\rho_{\epsilon}m\log(em)\).

For the heavy sum, condition further on the estimation-block category and treatment indicators.  By Assumption~\(\mathrm{ass:second\text{-}central\text{-}moment}\), its conditional outcome variance is bounded by
\[
 \frac{C}{m^2}\sum_{k\in\widehat{\mathcal H}}
 N_k^2\mathbf 1\{N_{0k}N_{1k}>0\}
 \left(\frac1{N_{1k}}+\frac1{N_{0k}}\right).
\]
Given \(N_k\), overlap and the elementary binomial reciprocal calculation used in Lemma~\(\mathrm{lem:continuous\text{-}occupancy\text{-}collision\text{-}upper}\) bound its expectation by \(C_{\epsilon}N_k/m^2\).  Summing \(\sum_kN_k=m\) yields \(C_{\epsilon}/m\).  After removing the missing-arm terms, the conditional mean is the empirical average of \((\vartheta_X-t)\mathbf 1\{X\in\widehat{\mathcal H}\}\), whose variance is at most \(4/m\).  By \((5)\), mean normalization, and \(\sum_kp_k=1\), the missing-arm remainder has second moment at most \(C_{\epsilon}n^{-8}\).  Therefore
\[
 E\!\left[\left.
 \left\{
 \sum_{k\in\widehat{\mathcal H}}\widehat h_k(t)
 -\sum_{k\in\widehat{\mathcal H}}p_k(\vartheta_k-t)
 \right\}^2\right|t,I_{\mathrm{cls}},\mathcal A\right]
 \le C_{\epsilon}/n+C_{\epsilon}n^{-8}.
\tag{13}
\]

Let \(Z(t)\) be the normalized expression before final clipping.  Combining \((9)\)--\((13)\) and using \((x+y+z)^2\le3(x^2+y^2+z^2)\) gives, with \(b:=n^{-1}+d/n^2\),
\[
 E[\{Z(t)-\vartheta\}^2\mid t,I_{\mathrm{cls}},\mathcal A]
 \le C_{\epsilon}\left[
 b+u\{\sigma^2+(t-\vartheta)^2\}
 \right].
\tag{14}
\]
Projection onto \([-1,1]\) cannot increase loss from \(\vartheta\in[-1,1]\), and the projected squared loss is at most four on \(\mathcal A^c\).  Integrating \((14)\), then applying \((2)\), \((4)\), and \(u\le\kappa_{\epsilon}<1\), yields
\[
\begin{aligned}
 E[(T_n^{\mathrm{ctr}}/M-\vartheta)^2]
 &\le C_{\epsilon}\{b+u\sigma^2+uE(t-\vartheta)^2+n^{-8}\}\\
 &\le C_{\epsilon}\{b+u\sigma^2+u(b+\sigma^2)\}\\
 &\le C_{\epsilon}(b+\sigma^2u).
\end{aligned}
\tag{15}
\]
Multiplication by \(M^2\) proves the risk claim.  Every statistic in the construction has a displayed zero fallback and final clipping maps it into \([-M,M]\), so \(T_n^{\mathrm{ctr}}\in\mathcal T_{n,M}\).  The only outcome-tail input is the conditional variance bound used in the centered covariance lemma and in \((13)\).